% =====================================================================
%  3. PROJEKT STACJI BAZOWEJ  (wersja 2026-10-29)
%  Objętość: ok. 9 stron. Model (35_Model.tex) dołączony na końcu.
% =====================================================================
\chapter{Projekt stacji bazowej}
\label{chap:projekt}

% 2-3 zdania: co jest w tym rozdziale.

% ---------------------------------------------------------------------
\section{Wymagania}
\label{sec:wymagania}

\subsection{Wymagania funkcjonalne}
\label{sec:wymagania_funkcjonalne}
% Odbiór ramek ADS-B, odczyt pozycji i czasu GPS, dekodowanie pozycji
% w czasie rzeczywistym, publikacja jako komunikaty ROS 2, rejestracja,
% podgląd stanu na wyświetlaczu.

\subsection{Wymagania pozafunkcjonalne}
\label{sec:wymagania_pozafunkcjonalne}
% Wielogodzinna praca z powerbanka, przenośność, start bez obsługi
% (autostart), odporność na nagłą utratę zasilania, praca bez internetu.

\subsection{Wymagania dotyczące danych referencyjnych}
\label{sec:wymagania_referencja}
% Znaczniki czasu UTC o znanej dokładności, znana pozycja stacji,
% bezstratny zapis surowych danych, jednoznaczne pozycje samolotów,
% kierunek do samolotu (azymut, elewacja).

% ---------------------------------------------------------------------
\section{Architektura systemu}
\label{sec:architektura}
% (ok. 2 str.) Podział na warstwę sprzętową i programową; podział
% demodulacja (dump1090) / dekodowanie (decoder_node); diagram przepływu
% danych; scenariusz pracy w terenie (włączenie -> nagrywanie ->
% wyłączenie przyciskiem -> skopiowanie danych).
% ======================================================================
% rys_sprzet.tex – schemat blokowy warstwy sprzętowej stacji bazowej
% Wstawić w rozdziale 3 (Architektura) poleceniem: \input{tex/rys_sprzet}
% Wymaga w preambule: tikz + biblioteki positioning, arrows.meta, calc,
%                     decorations.pathmorphing; siunitx; kolor akcent.
% Kod TikZ przygotowano z wykorzystaniem Claude (Anthropic) – zob. deklarację AI.
% ======================================================================

Warstwę sprzętową stacji bazowej przedstawiono na rysunku~\ref{rys:sprzet}. Centralnym elementem jest komputer Raspberry Pi~4B, do którego przez port USB dołączono odbiornik SDR, a przez interfejs UART moduł odbiornika GPS. Wyświetlacz TFT, dioda stanu i~przycisk zasilania są połączone bezpośrednio z~wyprowadzeniami GPIO.

\begin{figure}[!htb]
\centering
\begin{tikzpicture}[
    font=\footnotesize,
    >={Stealth[length=2mm]},
    blok/.style={draw, thick, fill=akcent!10, align=center,
                 minimum height=9mm, inner sep=3pt, text width=#1},
    blok/.default=27mm,
    zewn/.style={align=center, text width=34mm, font=\footnotesize},
    dane/.style={->, thick},
    zasil/.style={->, line width=2pt, draw=black!50},
    opcja/.style={<->, thick, dashed},
    radio/.style={->, thick, decorate,
                  decoration={snake, amplitude=0.6mm, segment length=2.5mm,
                              post length=2mm}},
    etyk/.style={font=\scriptsize, align=center, fill=white, inner sep=1pt},
]
% ---------------- Raspberry Pi (środek) ----------------
\node[draw, very thick, fill=akcent!20, align=center,
      minimum width=44mm, minimum height=44mm] (rpi) at (0,0)
     {\textbf{Raspberry Pi 4B}\\[2mm]
      \scriptsize Ubuntu, Docker,\\ \scriptsize ROS~2 Jazzy,\\
      \scriptsize dump1090};

% ---------------- lewa strona: odbiór radiowy ----------------
\node[blok] (sdr) at (-6.3, 1.3) {Odbiornik SDR\\RTL-SDR Blog V3};
\node[blok] (gps) at (-6.3,-1.3) {Moduł odbiornika GPS};
\node[zewn] (adsb) at (-6.3, 3.6)
     {statki powietrzne\\(ADS-B, \SI{1090}{\mega\hertz})};
\node[zewn] (sat) at (-6.3,-3.6)
     {satelity GPS\\(L1, \SI{1575.42}{\mega\hertz})};
\draw[radio] (adsb) -- node[etyk, right=1mm, fill=none]
     {antena dipolowa $\lambda/4$} (sdr);
\draw[radio] (sat) -- node[etyk, right=1mm, fill=none]
     {antena GPS} (gps);
\draw[dane] (sdr.east) -- node[etyk, above=1pt, fill=none] {USB~2.0}
     (sdr.east -| rpi.west);
\draw[dane] (gps.east) -- node[etyk, above=1pt, fill=none]
     {UART} node[etyk, below=1pt, fill=none] {\SI{9600}{bit/s}}
     (gps.east -| rpi.west);

% ---------------- prawa strona: interfejs użytkownika ----------------
\node[blok] (tft) at (6.3, 1.5) {Wyświetlacz TFT\\ILI9341, 240$\times$320};
\node[blok] (led) at (6.3, 0.0) {Dioda LED stanu};
\node[blok] (btn) at (6.3,-1.5) {Przycisk zasilania};
\draw[dane] (tft.west -| rpi.east) -- node[etyk, above=1pt, fill=none] {SPI0}
     node[etyk, below=1pt, fill=none] {GPIO24/25} (tft.west);
\draw[dane] (led.west -| rpi.east) -- node[etyk, above=1pt, fill=none]
     {GPIO17} (led.west);
\draw[dane] (btn.west) -- node[etyk, above=1pt, fill=none] {GPIO3}
     (btn.west -| rpi.east);

% ---------------- dół: zasilanie i pamięć ----------------
\node[blok=31mm] (pb)  at (-1.9,-4.0) {Powerbank\\LYNXER RPP-272};
\node[blok=31mm] (usb) at ( 1.9,-4.0) {Pamięć USB \SI{64}{GB}\\(system, pliki MCAP)};
\draw[zasil] (pb.north) -- node[etyk, left=1pt, fill=none]
     {USB-C, \SI{5}{V}} (pb.north |- rpi.south);
\draw[<->, thick] (usb.north) -- node[etyk, right=1pt, fill=none] {USB}
     (usb.north |- rpi.south);

% ---------------- góra: opcjonalny komputer ----------------
\node[blok=36mm, fill=white, dashed] (pc) at (0, 3.9)
     {Komputer do analizy\\(opcjonalnie)};
\draw[opcja] (pc.south) -- node[etyk, right=1pt, fill=none]
     {Wi-Fi, DDS / SSH} (rpi.north);
\end{tikzpicture}
\caption{Schemat blokowy warstwy sprzętowej stacji bazowej. Strzałki ciągłe
oznaczają kierunek przepływu danych, gruba szara strzałka -- zasilanie,
strzałki faliste -- sygnały radiowe, linia przerywana -- połączenie
opcjonalne. Przy strzałkach podano interfejsy i numery wyprowadzeń GPIO
(numeracja BCM).}
\label{rys:sprzet}
\end{figure}
\newpage
% ---------------------------------------------------------------------
\section{Tematy i~komunikaty}
\label{sec:tematy}
% (ok. 1 str.) Lista tematów (tabela) + definicja komunikatu AircraftState
% (pola, jednostki SI, NaN = brak wartości, stałe position_quality),
% uzasadnienie własnego typu zamiast JSON w std_msgs/String.

% ======================================================================
% rys_graf_ros.tex – graf obliczeniowy ROS 2 stacji bazowej
% (węzły, tematy, typy komunikatów, źródła danych w systemie gospodarza)
% Wstawić w rozdziale 3 (Architektura) poleceniem: \input{tex/rys_graf_ros}
% Wymaga w preambule: tikz + biblioteki positioning, arrows.meta, fit,
%                     backgrounds, calc; kolor akcent.
% Nazwy węzłów i tematów zgodne z repozytorium (commit 33bb6ff).
% Kod TikZ przygotowano z wykorzystaniem Claude (Anthropic) – zob. deklarację AI.
% ======================================================================
\begin{figure}[!htb]
\centering
\begin{tikzpicture}[
    font=\footnotesize,
    >={Stealth[length=1.8mm]},
    wezel/.style={draw, thick, rounded corners=3mm, fill=akcent!20,
                  align=center, minimum height=8mm, text width=27mm,
                  font=\scriptsize\ttfamily},
    temat/.style={draw, thick, fill=white, align=center,
                  minimum height=8mm, text width=32mm, inner sep=2pt},
    host/.style={draw, thick, fill=black!7, align=center,
                 minimum height=8mm, text width=23mm, font=\scriptsize},
    pol/.style={->, thick},
    etyk/.style={font=\scriptsize, align=center, inner sep=1pt},
]
% makro: nazwa tematu (czcionka maszynowa) + typ komunikatu (kursywa)
\newcommand{\tm}[2]{{\scriptsize\ttfamily #1}\\[-0.6mm]{\scriptsize\itshape #2}}

% ---------------- system gospodarza (źródła danych) ----------------
\node[host] (d1090) at (0, 3.65)
     {\textbf{dump1090}\\(demodulacja)\\TCP 30003: SBS\\TCP 30002: AVR};
\node[host] (uart)  at (0, 1.25) {\texttt{/dev/ttyAMA0}\\(UART, GPS)};
\node[host] (hwmon) at (0,-1.15) {\texttt{/sys/class/hwmon}\\(czujniki RPi)};
\draw[pol] ($(d1090.north)+(0,0.6)$) node[above, etyk] {z RTL-SDR (USB)}
     -- (d1090.north);

% ---------------- węzły publikujące ----------------
\node[wezel] (sdr)  at (3.5, 3.65) {sdr\_radar\_node};
\node[wezel] (gpsn) at (3.5, 1.25) {gps\_node};
\node[wezel] (pwr)  at (3.5,-1.15) {power\_monitor\_node};
\draw[pol] ([yshift=2mm]d1090.east)  -- ([yshift=2mm]sdr.west);
\draw[pol] ([yshift=-2mm]d1090.east) -- ([yshift=-2mm]sdr.west);
\draw[pol] (uart)  -- (gpsn);
\draw[pol] (hwmon) -- (pwr);

% ---------------- tematy ----------------
\node[temat] (raw)    at (6.95, 4.25) {\tm{/aircraft/raw}{std\_msgs/String}};
\node[temat] (frames) at (6.95, 3.05) {\tm{/aircraft/frames}{std\_msgs/String}};
\node[temat] (fix)    at (6.95, 1.85) {\tm{/gps/fix}{sensor\_msgs/NavSatFix}};
\node[temat] (tref)   at (6.95, 0.65) {\tm{/gps/time\_ref}{sensor\_msgs/TimeReference}};
\node[temat] (uv)     at (6.95,-0.55) {\tm{/system/undervoltage}{std\_msgs/Bool}};
\node[temat] (temp)   at (6.95,-1.75) {\tm{/system/cpu\_temp}{sensor\_msgs/Temperature}};

\draw[pol] (sdr.east)  -- (raw.west);
\draw[pol] (sdr.east)  -- (frames.west);
\draw[pol] (gpsn.east) -- (fix.west);
\draw[pol] (gpsn.east) -- (tref.west);
\draw[pol] (pwr.east)  -- (uv.west);
\draw[pol] (pwr.east)  -- (temp.west);

% ---------------- przetwarzanie i prezentacja ----------------
\node[wezel] (dec) at (10.45, 2.45) {decoder\_node};
\node[temat] (state) at (10.45, 0.75)
     {\tm{/aircraft/state}{aircraft\_interfaces/\\[-0.6mm]AircraftState}};
\node[wezel, text width=19mm] (disp) at (13.65,-0.55) {display\_node};
\node[etyk, below=1mm of disp.south east, anchor=north east] {$\rightarrow$ TFT (SPI)};

\draw[pol] (frames.east) -- ([yshift=1.5mm]dec.west);
\draw[pol] (fix.east)    -- ([yshift=-1.5mm]dec.west);
\draw[pol] (dec.south)   -- (state.north);
\draw[pol] (fix.east)   -| ([xshift=3mm]disp.north);
\draw[pol] (state.east) -| ([xshift=-3mm]disp.north);
\draw[pol] (uv.east)    -- (disp.west);
\draw[pol] (temp.east)  -| ([xshift=-6mm]disp.south);

% ---------------- moduł adsb_decoder (nie jest węzłem) ----------------
\node[draw, thick, dashed, fill=white, align=center, text width=29mm,
      minimum height=8mm, inner sep=2pt] (mod) at (10.95, 4.25)
     {{\scriptsize\ttfamily adsb\_decoder.py}\\[-0.6mm]
      {\scriptsize\itshape moduł Pythona (nie węzeł)}};
\draw[->, dashed, thick, >={Straight Barb[length=1.8mm]}]
     (mod.south |- dec.north) -- node[etyk, right=1mm] {«import»} (mod.south);

% ---------------- rejestracja rosbag2 ----------------
% ramka w kształcie litery L: kolumna tematów + temat /aircraft/state
% (bez węzła decoder_node, który nie jest tematem)
\coordinate (gL) at ($(raw.west)+(-2.5mm,0)$);
\coordinate (gR) at ($(raw.east)+(2.5mm,0)$);
\coordinate (gT) at ($(raw.north)+(0,2.5mm)$);
\coordinate (gB) at ($(temp.south)+(0,-2.5mm)$);
\coordinate (sT) at ($(state.north)+(0,2mm)$);
\coordinate (sB) at ($(state.south)+(0,-2mm)$);
\coordinate (sR) at ($(state.east)+(2mm,0)$);
\draw[dashed, rounded corners=2mm]
     (gL |- gT) -- (gR |- gT) -- (gR |- sT) -- (sR |- sT) -- (sR |- sB)
     -- (gR |- sB) -- (gR |- gB) -- (gL |- gB) -- cycle;
\node[wezel, text width=30mm] (bag) at (6.95,-3.4) {rosbag2\_recorder};
\node[etyk, right=2mm of bag, align=left] (mcap)
     {$\rightarrow$ pliki \texttt{.mcap}\\\phantom{$\rightarrow$} (co \SI{10}{min})};
\draw[->, line width=1.4pt] (gB -| bag.north) -- node[etyk, left=1mm]
     {wszystkie tematy z~ramki} (bag.north);

% ---------------- granice: Ubuntu (poza kontenerem) i kontener ----------------
\begin{scope}[on background layer]
  \node[draw, thick, densely dotted, rounded corners=2mm,
        fit=(sdr)(pwr)(raw)(temp)(state)(disp)(mod)(bag)(mcap), inner sep=3mm,
        label={[etyk, anchor=south east]north east:kontener Docker (ROS~2 Jazzy)}]
        (kont) {};
  \node[draw, thick, densely dotted, rounded corners=2mm,
        fit=(d1090)(hwmon), inner sep=1.5mm,
        label={[etyk, anchor=north, text width=25mm]south:Ubuntu na RPi\\(poza kontenerem)}] (gosp) {};
\end{scope}
\end{tikzpicture}
\caption{Graf obliczeniowy ROS~2 stacji bazowej. Prostokąty zaokrąglone
oznaczają węzły, prostokąty o~ostrych narożnikach -- tematy wraz z~typem
komunikatu (kursywa), szare prostokąty -- źródła danych w~systemie Ubuntu
uruchomionym bezpośrednio na Raspberry Pi, poza kontenerem. Strzałka ciągła
prowadzi od wydawcy do subskrybenta. Prostokąt przerywany oznacza moduł Pythona
\texttt{adsb\_decoder.py}, który nie jest węzłem ROS~2, a~przerywana strzałka
«import» -- jego wykorzystanie przez węzeł \texttt{decoder\_node}. Ramka
przerywana obejmuje tematy rejestrowane przez rosbag2, ramki kropkowane --
elementy uruchamiane w~kontenerze Docker i~poza nim.}
\label{rys:graf_ros}
\end{figure}

Przepływ danych w~środowisku ROS~2 pokazano na rysunku~\ref{rys:graf_ros}. Zgodnie z~typową praktyką w~ROS~2 każdy węzeł ma jedno, wyraźnie określone zadanie, a~dane między węzłami przekazywane są wyłącznie przez tematy. Węzeł \texttt{sdr\_radar\_node} odbiera od programu dump1090 zdekodowane wiadomości SBS oraz surowe ramki w~formacie AVR i~publikuje je w~dwóch osobnych tematach. Węzeł \texttt{gps\_node} odczytuje zdania NMEA z~odbiornika GPS, a~następnie publikuje pozycję stacji bazowej oraz czas GPS, każde w~osobnym temacie. Węzeł \texttt{decoder\_node} subskrybuje temat z~ramkami i~temat z~pozycją stacji, wyznacza na ich podstawie stan statków powietrznych względem stacji bazowej i~publikuje go w~temacie \texttt{/aircraft/state}. Do dekodowania ramek wykorzystuje moduł \texttt{adsb\_decoder.py}, który nie jest węzłem. Węzeł \texttt{power\_monitor\_node} odczytuje z~systemu temperaturę procesora oraz informację o~zbyt niskim napięciu zasilania i~publikuje je w~tematach \texttt{/system/cpu\_temp} i~\texttt{/system/undervoltage}. Węzeł \texttt{display\_node} subskrybuje te dwa tematy, a~także temat ze stanem statków powietrznych i~temat z~pozycją stacji, po czym wyświetla dane na ekranie TFT za pomocą biblioteki luma.lcd. Wszystkie tematy są zapisywane do plików MCAP przez rejestrator rosbag2.

\begin{table}[!htb]
  \centering
  \caption{Tematy ROS~2 publikowane przez stację bazową}
  \label{tab:tematy}
  \footnotesize
  \begin{tabular}{>{\raggedright}p{3.3cm} >{\raggedright}p{4.0cm} p{1.7cm} >{\raggedright\arraybackslash}p{4.9cm}}
    \toprule
    Temat & Typ komunikatu & Częstotl.\ [Hz] & Zawartość \\
    \midrule
    \texttt{/aircraft/raw}        & \texttt{std\_msgs/String}         & 0--50   & linia SBS z~dekodera dump1090 \\
    \texttt{/aircraft/frames}     & \texttt{std\_msgs/String}         & 30--70  & surowa ramka Mode~S w~formacie AVR \\
    \texttt{/aircraft/state}      & \texttt{aircraft\_interfaces/ AircraftState} & zmienna & stan statku powietrznego i~geometria względem stacji \\
    \texttt{/gps/fix}             & \texttt{sensor\_msgs/NavSatFix}   & 1       & pozycja stacji \\
    \texttt{/gps/time\_ref}       & \texttt{sensor\_msgs/ TimeReference} & 1    & czas UTC z~odbiornika GPS \\
    \texttt{/system/ undervoltage} & \texttt{std\_msgs/Bool}           & 1       & sygnalizacja zbyt niskiego napięcia zasilania \\
    \texttt{/system/cpu\_temp}    & \texttt{sensor\_msgs/ Temperature} & 1       & temperatura procesora \\
    \bottomrule
  \end{tabular}
\end{table}

% ---------------------------------------------------------------------
%  3.4. Model wyznaczania danych referencyjnych (osobny plik, zaczyna
%  się od \section). Umieszczony PO architekturze i tematach, bo model
%  korzysta z pojęć tam wprowadzonych (/gps/time_ref, znaczniki czasu).
% ---------------------------------------------------------------------
% ============================================================
%  35_Model.tex  --  Model opisowy i matematyczny wyznaczania
%                    danych referencyjnych (wersja uproszczona)
%  Wstawiany do rozdzialu 3 poleceniem:  % ============================================================
%  35_Model.tex  --  Model opisowy i matematyczny wyznaczania
%                    danych referencyjnych (wersja uproszczona)
%  Wstawiany do rozdzialu 3 poleceniem:  % ============================================================
%  35_Model.tex  --  Model opisowy i matematyczny wyznaczania
%                    danych referencyjnych (wersja uproszczona)
%  Wstawiany do rozdzialu 3 poleceniem:  \input{tex/35_Model}
%  Wymaga w preambule: amsmath, siunitx, enumitem
% ============================================================

% Objaśnienia symboli pod wzorem ("gdzie:"): symbol w kolumnie od
% wcięcia akapitu, myślniki zawsze w jednej pionowej linii
\newenvironment{objasnienia}{%
  \noindent gdzie:
  \begin{description}[font=\normalfont, labelindent=\parindent,
    labelwidth=1.2cm, labelsep=0.2cm, leftmargin=!, align=left,
    itemsep=0pt, topsep=2pt]}{\end{description}}

\section{Model wyznaczania danych referencyjnych}
\label{sec:model}

Zarejestrowane przez stację dane mają posłużyć jako dane referencyjne, czyli
mają pozwolić stwierdzić, gdzie znajdował się dany statek powietrzny
w określonej chwili czasu. % referencja do podrozdziału.
W tym celu potrzebny jest model, który łączy trzy źródła informacji:
komunikaty ADS-B odbierane drogą radiową, pozycję stacji z odbiornika GPS oraz
zegar systemowy komputera, którym oznaczane są wszystkie zapisane wiadomości.
W niniejszym podrozdziale model ten przedstawiono najpierw w postaci słownej
(model opisowy), a następnie w postaci wzorów (model matematyczny). Na końcu
omówiono przyjęte uproszczenia i podano przykład obliczeniowy.

\subsection{Model opisowy}
\label{sec:model_opisowy}

Statek powietrzny wyposażony w transponder ADS-B kilka razy na sekundę
rozgłasza na częstotliwości \SI{1090}{MHz} komunikaty zawierające m.in.\ jego
unikalny adres ICAO, pozycję (szerokość i długość geograficzną) wyznaczoną
przez pokładowy odbiornik GNSS, wysokość barometryczną oraz prędkość
względem ziemi~\cite{sun2021}. Stacja odbiera część tych komunikatów --
niektóre giną z powodu zakłóceń, zbyt słabego sygnału lub przesłonięcia
anteny. Odebrane komunikaty są zapisywane w postaci tekstowej (format SBS)
razem ze znacznikiem czasu nadawanym przez komputer w chwili odbioru.

Komunikat ADS-B mówi, gdzie samolot jest, ale nie mówi, gdzie jest
względem stacji. Do tego potrzebna jest pozycja samej stacji, którą
podaje odbiornik GPS w komunikacie NMEA \texttt{GGA}~\cite{nmea0183}.
Znając obie pozycje, można obliczyć odległość od stacji do samolotu.

Drugim problemem jest czas. Komputer Raspberry Pi nie ma podtrzymywanego
bateryjnie zegara, a w terenie nie ma dostępu do internetu, więc jego zegar
może się mylić o nieznaną wartość. Ponieważ wszystkie zapisane wiadomości są
oznaczone właśnie tym zegarem, błąd zegara przenosi się na wszystkie dane.
Odbiornik GPS co sekundę przesyła jednak dokładny czas UTC (komunikat
\texttt{RMC}). Komputer zapisuje go razem z własnym wskazaniem zegara, dzięki
czemu po zakończeniu pomiarów można zmierzyć, o ile zegar się mylił,
i poprawić czas wszystkich wiadomości.

Granice modelowanego układu wyznacza sama stacja (antena, odbiornik RTL-SDR,
odbiornik GPS i komputer). Samoloty i satelity GPS należą do otoczenia,
a punktami styku są: łącze radiowe \SI{1090}{MHz}, łącze szeregowe
odbiornika GPS oraz plik zapisu rosbag2. Model składa się z dwóch części:
\begin{enumerate}[nosep]
  \item modelu czasu -- poprawienie znaczników czasu zegara systemowego
        na czas UTC (podrozdział~\ref{sec:model_czasu}),
  \item modelu odległości -- obliczenie odległości od stacji do samolotu
        (podrozdział~\ref{sec:model_odleglosc}).
\end{enumerate}

\subsection{Model czasu}
\label{sec:model_czasu}

Niech $t$ oznacza wskazanie zegara systemowego komputera, a $\tau$ --
odpowiadający mu rzeczywisty czas UTC. Zależność między nimi zapisano jako
\begin{equation}
  \tau = t + \delta,
  \label{eq:model_zegara}
\end{equation}
gdzie $\delta$ jest przesunięciem (błędem) zegara systemowego w sekundach.
Przyjęto, że w trakcie jednej sesji pomiarowej przesunięcie $\delta$ jest
stałe.

Węzeł obsługi GPS publikuje w temacie \texttt{/gps/time\_ref} pary
$(t_k, \tau_k)$, $k = 1, 2, \ldots, K$, gdzie $\tau_k$ jest czasem UTC
odczytanym z komunikatu \texttt{RMC}, a $t_k$ -- wskazaniem zegara
systemowego w chwili odczytu tego komunikatu. Każda para daje jeden pomiar
przesunięcia:
\begin{equation}
  \delta_k = \tau_k - t_k, \qquad k = 1, 2, \ldots, K.
  \label{eq:offset_pomiar}
\end{equation}
Pojedyncze pomiary różnią się nieco od siebie, dlatego jako wynik przyjmuje
się ich średnią arytmetyczną, a rozrzut pomiarów opisuje się odchyleniem
standardowym:
\begin{equation}
  \hat{\delta} = \frac{1}{K} \sum_{k=1}^{K} \delta_k,
  \qquad
  \sigma_{\delta} = \sqrt{\frac{1}{K-1} \sum_{k=1}^{K}
                    \bigl(\delta_k - \hat{\delta}\bigr)^2},
  \label{eq:offset_srednia}
\end{equation}
\begin{objasnienia}
  \item[$\hat{\delta}$] -- wyznaczone przesunięcie zegara w sekundach,
  \item[$\sigma_{\delta}$] -- odchylenie standardowe pomiarów przesunięcia w sekundach,
  \item[$K$] -- liczba pomiarów (przy jednym pomiarze na sekundę sesja trwająca \SI{10}{min} daje ok.~600 pomiarów).
\end{objasnienia}
Poprawiony czas UTC dowolnej zapisanej wiadomości o znaczniku $t$ wynosi
\begin{equation}
  \hat{\tau} = t + \hat{\delta}.
  \label{eq:korekta_czasu}
\end{equation}

Pomiar \eqref{eq:offset_pomiar} nie jest idealny, ponieważ znacznik $t_k$
jest nadawany dopiero po przesłaniu komunikatu łączem szeregowym i odczytaniu
go przez program. Przy prędkości \num{9600}~bit/s przesłanie jednego znaku
(10~bitów) trwa ok.\ \SI{1.04}{ms}, a komunikat NMEA ma do 82 znaków, więc
sama transmisja zajmuje do ok.\ \SI{85}{ms}. Program odczytuje port co
\SI{100}{ms}, co dodaje średnio kolejne \SI{50}{ms}. Wyznaczone przesunięcie
$\hat{\delta}$ jest więc obarczone błędem rzędu \SI{0.1}{s}, co -- jak
pokazano w podrozdziale~\ref{sec:model_przyklad} -- jest wartością
akceptowalną.

\subsection{Model odległości}
\label{sec:model_odleglosc}

\paragraph{Dane wejściowe.}
Z komunikatu SBS typu \texttt{MSG,3} odczytuje się szerokość $\varphi_i$
i długość $\lambda_i$ geograficzną samolotu (w stopniach) oraz jego wysokość
w stopach. Z tematu \texttt{/gps/fix} odczytuje się szerokość
$\varphi_{\mathrm{s}}$, długość $\lambda_{\mathrm{s}}$ i wysokość
$h_{\mathrm{s}}$ stacji (w metrach). Indeks $i$ oznacza samolot, a indeks
$\mathrm{s}$ -- stację.

\paragraph{Przeliczenie jednostek.}
W lotnictwie wysokość podaje się w stopach, a prędkość w węzłach. Na
jednostki układu SI przelicza się je wzorami
\begin{equation}
  h_i = \num{0.3048}\, h_{\mathrm{ft}}, \qquad
  v = \num{0.514444}\, v_{\mathrm{kt}},
  \label{eq:jednostki}
\end{equation}
\begin{objasnienia}
  \item[$h_{\mathrm{ft}}$] -- wysokość samolotu w stopach,
  \item[$h_i$] -- wysokość samolotu w metrach,
  \item[$v_{\mathrm{kt}}$] -- prędkość w węzłach,
  \item[$v$] -- prędkość w metrach na sekundę.
\end{objasnienia}

\paragraph{Odległość po powierzchni Ziemi.}
Ziemię przybliżono kulą o promieniu $R_{\mathrm{Z}} = \SI{6371}{km}$.
Odległość między stacją a punktem na powierzchni Ziemi leżącym dokładnie pod
samolotem oblicza się wzorem haversine:
\begin{equation}
  s = 2 R_{\mathrm{Z}} \arcsin \sqrt{
        \sin^2\frac{\Delta\varphi}{2}
        + \cos\varphi_{\mathrm{s}} \cos\varphi_i
          \sin^2\frac{\Delta\lambda}{2}},
  \label{eq:haversine}
\end{equation}
\begin{objasnienia}
  \item[$s$] -- odległość mierzona po powierzchni Ziemi w metrach,
  \item[$\Delta\varphi$] -- różnica szerokości geograficznych, $\Delta\varphi = \varphi_i - \varphi_{\mathrm{s}}$, w radianach,
  \item[$\Delta\lambda$] -- różnica długości geograficznych, $\Delta\lambda = \lambda_i - \lambda_{\mathrm{s}}$, w radianach,
  \item[$R_{\mathrm{Z}}$] -- średni promień Ziemi w metrach.
\end{objasnienia}
Kąty podawane w stopniach należy przed użyciem we wzorze
\eqref{eq:haversine} zamienić na radiany, mnożąc je przez $\pi/180$.

\paragraph{Odległość skośna.}
Samolot znajduje się na wysokości $h_i$, a stacja na wysokości
$h_{\mathrm{s}}$. Odległość w linii prostej od stacji do samolotu
(odległość skośną) przybliża się, traktując odległość $s$ i różnicę wysokości
jako przyprostokątne trójkąta prostokątnego:
\begin{equation}
  \rho = \sqrt{s^2 + \Delta h^2}, \qquad \Delta h = h_i - h_{\mathrm{s}},
  \label{eq:odleglosc_skosna}
\end{equation}
gdzie $\rho$ jest odległością skośną, a $\Delta h$ -- różnicą wysokości
samolotu i stacji, obie w metrach. Wzory \eqref{eq:jednostki}%
--\eqref{eq:odleglosc_skosna} są zaimplementowane w węźle wyświetlacza
stacji, który na ich podstawie pokazuje odległość do każdego samolotu.

\subsection{Uproszczenia modelu}
\label{sec:model_uproszczenia}

W modelu przyjęto następujące uproszczenia:
\begin{itemize}[nosep]
  \item Ziemia jest kulą, a nie elipsoidą -- powoduje to błąd odległości
        rzędu kilku dziesiątych procenta,
  \item we wzorze \eqref{eq:odleglosc_skosna} pominięto krzywiznę Ziemi;
        błąd rośnie wraz z odległością samolotu od stacji,
  \item wysokość barometryczna z ADS-B jest traktowana jak wysokość nad
        poziomem morza; przy ciśnieniu innym niż standardowe
        (\SI{1013.25}{hPa}) różnica może sięgać kilkudziesięciu metrów,
  \item przesunięcie zegara $\delta$ jest stałe w trakcie sesji (pominięto
        powolny dryf zegara),
  \item stacja jest nieruchoma w trakcie sesji,
  \item pominięto opóźnienie między odebraniem komunikatu ADS-B a nadaniem mu
        znacznika czasu (pojedyncze milisekundy).
\end{itemize}
Uproszczenia te są wystarczające do podglądu na żywo i do analizy zasięgu
stacji. Przy dalszym rozwoju systemu model odległości można zastąpić
dokładniejszym przeliczeniem współrzędnych do lokalnego układu kartezjańskiego
związanego ze stacją, które uwzględnia kształt Ziemi i pozwala dodatkowo
wyznaczyć kierunek obserwacji samolotu.

\subsection{Przykład obliczeniowy}
\label{sec:model_przyklad}

Przyjęto stację w punkcie $\varphi_{\mathrm{s}} = \SI{50.0}{\degree}$,
$\lambda_{\mathrm{s}} = \SI{20.0}{\degree}$, $h_{\mathrm{s}} = \SI{250}{m}$
oraz samolot w punkcie $\varphi_i = \SI{50.3169}{\degree}$,
$\lambda_i = \SI{20.49795}{\degree}$, lecący na wysokości
$h_{\mathrm{ft}} = \num{32800}$~stóp.

Z wzoru \eqref{eq:jednostki} wysokość samolotu wynosi
\begin{equation}
  h_i = \num{0.3048} \cdot \num{32800} \approx \SI{9997.4}{m},
  \qquad
  \Delta h = \num{9997.4} - 250 = \SI{9747.4}{m}.
\end{equation}
Różnice współrzędnych wynoszą $\Delta\varphi = \SI{0.3169}{\degree}$
i $\Delta\lambda = \SI{0.49795}{\degree}$. Wyrażenie pod pierwiastkiem we
wzorze \eqref{eq:haversine} ma wartość ok.\ \num{1.540e-5}, skąd
\begin{equation}
  s \approx \SI{50.00}{km}, \qquad
  \rho = \sqrt{\num{50000}^2 + \num{9747.4}^2} \approx \SI{50.94}{km}.
\end{equation}

Wpływ błędu czasu pokazuje następujące oszacowanie: samolot lecący
z prędkością $v = \SI{230}{m/s}$ w czasie $\Delta\tau$ pokonuje drogę
\begin{equation}
  \Delta x = v \, \Delta\tau .
  \label{eq:blad_czasu}
\end{equation}
Jeżeli zegar stacji myli się o \SI{10}{s} i nie zostanie poprawiony wzorem
\eqref{eq:korekta_czasu}, pozycja przypisana danej chwili będzie błędna
o ok.\ \SI{2.3}{km}. Po korekcji błąd czasu spada do ok.\ \SI{0.1}{s},
a błąd pozycji -- do ok.\ \SI{23}{m}. Uzasadnia to rejestrowanie czasu GPS
razem z danymi ADS-B.

\subsection{Weryfikacja modelu}
\label{sec:model_weryfikacja}

Poprawność modelu odległości można sprawdzić na prostych przypadkach:
dla samolotu położonego dokładnie nad stacją powinno zachodzić $s = 0$
i $\rho = \Delta h$, a dla punktów odległych o \SI{1}{\degree} szerokości
geograficznej $s \approx \SI{111.2}{km}$. Poprawność modelu czasu można
sprawdzić, porównując $\hat{\delta}$ wyznaczone w sesji z dostępem do
internetu (gdy zegar jest synchronizowany przez NTP i $\hat{\delta}$ powinno
być bliskie zera) z sesją bez dostępu do sieci. Szczegółową ocenę na
podstawie zarejestrowanych danych przedstawiono w~rozdziale~\ref{chap:badania}.
%  Wymaga w preambule: amsmath, siunitx, enumitem
% ============================================================

% Objaśnienia symboli pod wzorem ("gdzie:"): symbol w kolumnie od
% wcięcia akapitu, myślniki zawsze w jednej pionowej linii
\newenvironment{objasnienia}{%
  \noindent gdzie:
  \begin{description}[font=\normalfont, labelindent=\parindent,
    labelwidth=1.2cm, labelsep=0.2cm, leftmargin=!, align=left,
    itemsep=0pt, topsep=2pt]}{\end{description}}

\section{Model wyznaczania danych referencyjnych}
\label{sec:model}

Zarejestrowane przez stację dane mają posłużyć jako dane referencyjne, czyli
mają pozwolić stwierdzić, gdzie znajdował się dany statek powietrzny
w określonej chwili czasu. % referencja do podrozdziału.
W tym celu potrzebny jest model, który łączy trzy źródła informacji:
komunikaty ADS-B odbierane drogą radiową, pozycję stacji z odbiornika GPS oraz
zegar systemowy komputera, którym oznaczane są wszystkie zapisane wiadomości.
W niniejszym podrozdziale model ten przedstawiono najpierw w postaci słownej
(model opisowy), a następnie w postaci wzorów (model matematyczny). Na końcu
omówiono przyjęte uproszczenia i podano przykład obliczeniowy.

\subsection{Model opisowy}
\label{sec:model_opisowy}

Statek powietrzny wyposażony w transponder ADS-B kilka razy na sekundę
rozgłasza na częstotliwości \SI{1090}{MHz} komunikaty zawierające m.in.\ jego
unikalny adres ICAO, pozycję (szerokość i długość geograficzną) wyznaczoną
przez pokładowy odbiornik GNSS, wysokość barometryczną oraz prędkość
względem ziemi~\cite{sun2021}. Stacja odbiera część tych komunikatów --
niektóre giną z powodu zakłóceń, zbyt słabego sygnału lub przesłonięcia
anteny. Odebrane komunikaty są zapisywane w postaci tekstowej (format SBS)
razem ze znacznikiem czasu nadawanym przez komputer w chwili odbioru.

Komunikat ADS-B mówi, gdzie samolot jest, ale nie mówi, gdzie jest
względem stacji. Do tego potrzebna jest pozycja samej stacji, którą
podaje odbiornik GPS w komunikacie NMEA \texttt{GGA}~\cite{nmea0183}.
Znając obie pozycje, można obliczyć odległość od stacji do samolotu.

Drugim problemem jest czas. Komputer Raspberry Pi nie ma podtrzymywanego
bateryjnie zegara, a w terenie nie ma dostępu do internetu, więc jego zegar
może się mylić o nieznaną wartość. Ponieważ wszystkie zapisane wiadomości są
oznaczone właśnie tym zegarem, błąd zegara przenosi się na wszystkie dane.
Odbiornik GPS co sekundę przesyła jednak dokładny czas UTC (komunikat
\texttt{RMC}). Komputer zapisuje go razem z własnym wskazaniem zegara, dzięki
czemu po zakończeniu pomiarów można zmierzyć, o ile zegar się mylił,
i poprawić czas wszystkich wiadomości.

Granice modelowanego układu wyznacza sama stacja (antena, odbiornik RTL-SDR,
odbiornik GPS i komputer). Samoloty i satelity GPS należą do otoczenia,
a punktami styku są: łącze radiowe \SI{1090}{MHz}, łącze szeregowe
odbiornika GPS oraz plik zapisu rosbag2. Model składa się z dwóch części:
\begin{enumerate}[nosep]
  \item modelu czasu -- poprawienie znaczników czasu zegara systemowego
        na czas UTC (podrozdział~\ref{sec:model_czasu}),
  \item modelu odległości -- obliczenie odległości od stacji do samolotu
        (podrozdział~\ref{sec:model_odleglosc}).
\end{enumerate}

\subsection{Model czasu}
\label{sec:model_czasu}

Niech $t$ oznacza wskazanie zegara systemowego komputera, a $\tau$ --
odpowiadający mu rzeczywisty czas UTC. Zależność między nimi zapisano jako
\begin{equation}
  \tau = t + \delta,
  \label{eq:model_zegara}
\end{equation}
gdzie $\delta$ jest przesunięciem (błędem) zegara systemowego w sekundach.
Przyjęto, że w trakcie jednej sesji pomiarowej przesunięcie $\delta$ jest
stałe.

Węzeł obsługi GPS publikuje w temacie \texttt{/gps/time\_ref} pary
$(t_k, \tau_k)$, $k = 1, 2, \ldots, K$, gdzie $\tau_k$ jest czasem UTC
odczytanym z komunikatu \texttt{RMC}, a $t_k$ -- wskazaniem zegara
systemowego w chwili odczytu tego komunikatu. Każda para daje jeden pomiar
przesunięcia:
\begin{equation}
  \delta_k = \tau_k - t_k, \qquad k = 1, 2, \ldots, K.
  \label{eq:offset_pomiar}
\end{equation}
Pojedyncze pomiary różnią się nieco od siebie, dlatego jako wynik przyjmuje
się ich średnią arytmetyczną, a rozrzut pomiarów opisuje się odchyleniem
standardowym:
\begin{equation}
  \hat{\delta} = \frac{1}{K} \sum_{k=1}^{K} \delta_k,
  \qquad
  \sigma_{\delta} = \sqrt{\frac{1}{K-1} \sum_{k=1}^{K}
                    \bigl(\delta_k - \hat{\delta}\bigr)^2},
  \label{eq:offset_srednia}
\end{equation}
\begin{objasnienia}
  \item[$\hat{\delta}$] -- wyznaczone przesunięcie zegara w sekundach,
  \item[$\sigma_{\delta}$] -- odchylenie standardowe pomiarów przesunięcia w sekundach,
  \item[$K$] -- liczba pomiarów (przy jednym pomiarze na sekundę sesja trwająca \SI{10}{min} daje ok.~600 pomiarów).
\end{objasnienia}
Poprawiony czas UTC dowolnej zapisanej wiadomości o znaczniku $t$ wynosi
\begin{equation}
  \hat{\tau} = t + \hat{\delta}.
  \label{eq:korekta_czasu}
\end{equation}

Pomiar \eqref{eq:offset_pomiar} nie jest idealny, ponieważ znacznik $t_k$
jest nadawany dopiero po przesłaniu komunikatu łączem szeregowym i odczytaniu
go przez program. Przy prędkości \num{9600}~bit/s przesłanie jednego znaku
(10~bitów) trwa ok.\ \SI{1.04}{ms}, a komunikat NMEA ma do 82 znaków, więc
sama transmisja zajmuje do ok.\ \SI{85}{ms}. Program odczytuje port co
\SI{100}{ms}, co dodaje średnio kolejne \SI{50}{ms}. Wyznaczone przesunięcie
$\hat{\delta}$ jest więc obarczone błędem rzędu \SI{0.1}{s}, co -- jak
pokazano w podrozdziale~\ref{sec:model_przyklad} -- jest wartością
akceptowalną.

\subsection{Model odległości}
\label{sec:model_odleglosc}

\paragraph{Dane wejściowe.}
Z komunikatu SBS typu \texttt{MSG,3} odczytuje się szerokość $\varphi_i$
i długość $\lambda_i$ geograficzną samolotu (w stopniach) oraz jego wysokość
w stopach. Z tematu \texttt{/gps/fix} odczytuje się szerokość
$\varphi_{\mathrm{s}}$, długość $\lambda_{\mathrm{s}}$ i wysokość
$h_{\mathrm{s}}$ stacji (w metrach). Indeks $i$ oznacza samolot, a indeks
$\mathrm{s}$ -- stację.

\paragraph{Przeliczenie jednostek.}
W lotnictwie wysokość podaje się w stopach, a prędkość w węzłach. Na
jednostki układu SI przelicza się je wzorami
\begin{equation}
  h_i = \num{0.3048}\, h_{\mathrm{ft}}, \qquad
  v = \num{0.514444}\, v_{\mathrm{kt}},
  \label{eq:jednostki}
\end{equation}
\begin{objasnienia}
  \item[$h_{\mathrm{ft}}$] -- wysokość samolotu w stopach,
  \item[$h_i$] -- wysokość samolotu w metrach,
  \item[$v_{\mathrm{kt}}$] -- prędkość w węzłach,
  \item[$v$] -- prędkość w metrach na sekundę.
\end{objasnienia}

\paragraph{Odległość po powierzchni Ziemi.}
Ziemię przybliżono kulą o promieniu $R_{\mathrm{Z}} = \SI{6371}{km}$.
Odległość między stacją a punktem na powierzchni Ziemi leżącym dokładnie pod
samolotem oblicza się wzorem haversine:
\begin{equation}
  s = 2 R_{\mathrm{Z}} \arcsin \sqrt{
        \sin^2\frac{\Delta\varphi}{2}
        + \cos\varphi_{\mathrm{s}} \cos\varphi_i
          \sin^2\frac{\Delta\lambda}{2}},
  \label{eq:haversine}
\end{equation}
\begin{objasnienia}
  \item[$s$] -- odległość mierzona po powierzchni Ziemi w metrach,
  \item[$\Delta\varphi$] -- różnica szerokości geograficznych, $\Delta\varphi = \varphi_i - \varphi_{\mathrm{s}}$, w radianach,
  \item[$\Delta\lambda$] -- różnica długości geograficznych, $\Delta\lambda = \lambda_i - \lambda_{\mathrm{s}}$, w radianach,
  \item[$R_{\mathrm{Z}}$] -- średni promień Ziemi w metrach.
\end{objasnienia}
Kąty podawane w stopniach należy przed użyciem we wzorze
\eqref{eq:haversine} zamienić na radiany, mnożąc je przez $\pi/180$.

\paragraph{Odległość skośna.}
Samolot znajduje się na wysokości $h_i$, a stacja na wysokości
$h_{\mathrm{s}}$. Odległość w linii prostej od stacji do samolotu
(odległość skośną) przybliża się, traktując odległość $s$ i różnicę wysokości
jako przyprostokątne trójkąta prostokątnego:
\begin{equation}
  \rho = \sqrt{s^2 + \Delta h^2}, \qquad \Delta h = h_i - h_{\mathrm{s}},
  \label{eq:odleglosc_skosna}
\end{equation}
gdzie $\rho$ jest odległością skośną, a $\Delta h$ -- różnicą wysokości
samolotu i stacji, obie w metrach. Wzory \eqref{eq:jednostki}%
--\eqref{eq:odleglosc_skosna} są zaimplementowane w węźle wyświetlacza
stacji, który na ich podstawie pokazuje odległość do każdego samolotu.

\subsection{Uproszczenia modelu}
\label{sec:model_uproszczenia}

W modelu przyjęto następujące uproszczenia:
\begin{itemize}[nosep]
  \item Ziemia jest kulą, a nie elipsoidą -- powoduje to błąd odległości
        rzędu kilku dziesiątych procenta,
  \item we wzorze \eqref{eq:odleglosc_skosna} pominięto krzywiznę Ziemi;
        błąd rośnie wraz z odległością samolotu od stacji,
  \item wysokość barometryczna z ADS-B jest traktowana jak wysokość nad
        poziomem morza; przy ciśnieniu innym niż standardowe
        (\SI{1013.25}{hPa}) różnica może sięgać kilkudziesięciu metrów,
  \item przesunięcie zegara $\delta$ jest stałe w trakcie sesji (pominięto
        powolny dryf zegara),
  \item stacja jest nieruchoma w trakcie sesji,
  \item pominięto opóźnienie między odebraniem komunikatu ADS-B a nadaniem mu
        znacznika czasu (pojedyncze milisekundy).
\end{itemize}
Uproszczenia te są wystarczające do podglądu na żywo i do analizy zasięgu
stacji. Przy dalszym rozwoju systemu model odległości można zastąpić
dokładniejszym przeliczeniem współrzędnych do lokalnego układu kartezjańskiego
związanego ze stacją, które uwzględnia kształt Ziemi i pozwala dodatkowo
wyznaczyć kierunek obserwacji samolotu.

\subsection{Przykład obliczeniowy}
\label{sec:model_przyklad}

Przyjęto stację w punkcie $\varphi_{\mathrm{s}} = \SI{50.0}{\degree}$,
$\lambda_{\mathrm{s}} = \SI{20.0}{\degree}$, $h_{\mathrm{s}} = \SI{250}{m}$
oraz samolot w punkcie $\varphi_i = \SI{50.3169}{\degree}$,
$\lambda_i = \SI{20.49795}{\degree}$, lecący na wysokości
$h_{\mathrm{ft}} = \num{32800}$~stóp.

Z wzoru \eqref{eq:jednostki} wysokość samolotu wynosi
\begin{equation}
  h_i = \num{0.3048} \cdot \num{32800} \approx \SI{9997.4}{m},
  \qquad
  \Delta h = \num{9997.4} - 250 = \SI{9747.4}{m}.
\end{equation}
Różnice współrzędnych wynoszą $\Delta\varphi = \SI{0.3169}{\degree}$
i $\Delta\lambda = \SI{0.49795}{\degree}$. Wyrażenie pod pierwiastkiem we
wzorze \eqref{eq:haversine} ma wartość ok.\ \num{1.540e-5}, skąd
\begin{equation}
  s \approx \SI{50.00}{km}, \qquad
  \rho = \sqrt{\num{50000}^2 + \num{9747.4}^2} \approx \SI{50.94}{km}.
\end{equation}

Wpływ błędu czasu pokazuje następujące oszacowanie: samolot lecący
z prędkością $v = \SI{230}{m/s}$ w czasie $\Delta\tau$ pokonuje drogę
\begin{equation}
  \Delta x = v \, \Delta\tau .
  \label{eq:blad_czasu}
\end{equation}
Jeżeli zegar stacji myli się o \SI{10}{s} i nie zostanie poprawiony wzorem
\eqref{eq:korekta_czasu}, pozycja przypisana danej chwili będzie błędna
o ok.\ \SI{2.3}{km}. Po korekcji błąd czasu spada do ok.\ \SI{0.1}{s},
a błąd pozycji -- do ok.\ \SI{23}{m}. Uzasadnia to rejestrowanie czasu GPS
razem z danymi ADS-B.

\subsection{Weryfikacja modelu}
\label{sec:model_weryfikacja}

Poprawność modelu odległości można sprawdzić na prostych przypadkach:
dla samolotu położonego dokładnie nad stacją powinno zachodzić $s = 0$
i $\rho = \Delta h$, a dla punktów odległych o \SI{1}{\degree} szerokości
geograficznej $s \approx \SI{111.2}{km}$. Poprawność modelu czasu można
sprawdzić, porównując $\hat{\delta}$ wyznaczone w sesji z dostępem do
internetu (gdy zegar jest synchronizowany przez NTP i $\hat{\delta}$ powinno
być bliskie zera) z sesją bez dostępu do sieci. Szczegółową ocenę na
podstawie zarejestrowanych danych przedstawiono w~rozdziale~\ref{chap:badania}.
%  Wymaga w preambule: amsmath, siunitx, enumitem
% ============================================================

% Objaśnienia symboli pod wzorem ("gdzie:"): symbol w kolumnie od
% wcięcia akapitu, myślniki zawsze w jednej pionowej linii
\newenvironment{objasnienia}{%
  \noindent gdzie:
  \begin{description}[font=\normalfont, labelindent=\parindent,
    labelwidth=1.2cm, labelsep=0.2cm, leftmargin=!, align=left,
    itemsep=0pt, topsep=2pt]}{\end{description}}

\section{Model wyznaczania danych referencyjnych}
\label{sec:model}

Zarejestrowane przez stację dane mają posłużyć jako dane referencyjne, czyli
mają pozwolić stwierdzić, gdzie znajdował się dany statek powietrzny
w określonej chwili czasu. % referencja do podrozdziału.
W tym celu potrzebny jest model, który łączy trzy źródła informacji:
komunikaty ADS-B odbierane drogą radiową, pozycję stacji z odbiornika GPS oraz
zegar systemowy komputera, którym oznaczane są wszystkie zapisane wiadomości.
W niniejszym podrozdziale model ten przedstawiono najpierw w postaci słownej
(model opisowy), a następnie w postaci wzorów (model matematyczny). Na końcu
omówiono przyjęte uproszczenia i podano przykład obliczeniowy.

\subsection{Model opisowy}
\label{sec:model_opisowy}

Statek powietrzny wyposażony w transponder ADS-B kilka razy na sekundę
rozgłasza na częstotliwości \SI{1090}{MHz} komunikaty zawierające m.in.\ jego
unikalny adres ICAO, pozycję (szerokość i długość geograficzną) wyznaczoną
przez pokładowy odbiornik GNSS, wysokość barometryczną oraz prędkość
względem ziemi~\cite{sun2021}. Stacja odbiera część tych komunikatów --
niektóre giną z powodu zakłóceń, zbyt słabego sygnału lub przesłonięcia
anteny. Odebrane komunikaty są zapisywane w postaci tekstowej (format SBS)
razem ze znacznikiem czasu nadawanym przez komputer w chwili odbioru.

Komunikat ADS-B mówi, gdzie samolot jest, ale nie mówi, gdzie jest
względem stacji. Do tego potrzebna jest pozycja samej stacji, którą
podaje odbiornik GPS w komunikacie NMEA \texttt{GGA}~\cite{nmea0183}.
Znając obie pozycje, można obliczyć odległość od stacji do samolotu.

Drugim problemem jest czas. Komputer Raspberry Pi nie ma podtrzymywanego
bateryjnie zegara, a w terenie nie ma dostępu do internetu, więc jego zegar
może się mylić o nieznaną wartość. Ponieważ wszystkie zapisane wiadomości są
oznaczone właśnie tym zegarem, błąd zegara przenosi się na wszystkie dane.
Odbiornik GPS co sekundę przesyła jednak dokładny czas UTC (komunikat
\texttt{RMC}). Komputer zapisuje go razem z własnym wskazaniem zegara, dzięki
czemu po zakończeniu pomiarów można zmierzyć, o ile zegar się mylił,
i poprawić czas wszystkich wiadomości.

Granice modelowanego układu wyznacza sama stacja (antena, odbiornik RTL-SDR,
odbiornik GPS i komputer). Samoloty i satelity GPS należą do otoczenia,
a punktami styku są: łącze radiowe \SI{1090}{MHz}, łącze szeregowe
odbiornika GPS oraz plik zapisu rosbag2. Model składa się z dwóch części:
\begin{enumerate}[nosep]
  \item modelu czasu -- poprawienie znaczników czasu zegara systemowego
        na czas UTC (podrozdział~\ref{sec:model_czasu}),
  \item modelu odległości -- obliczenie odległości od stacji do samolotu
        (podrozdział~\ref{sec:model_odleglosc}).
\end{enumerate}

\subsection{Model czasu}
\label{sec:model_czasu}

Niech $t$ oznacza wskazanie zegara systemowego komputera, a $\tau$ --
odpowiadający mu rzeczywisty czas UTC. Zależność między nimi zapisano jako
\begin{equation}
  \tau = t + \delta,
  \label{eq:model_zegara}
\end{equation}
gdzie $\delta$ jest przesunięciem (błędem) zegara systemowego w sekundach.
Przyjęto, że w trakcie jednej sesji pomiarowej przesunięcie $\delta$ jest
stałe.

Węzeł obsługi GPS publikuje w temacie \texttt{/gps/time\_ref} pary
$(t_k, \tau_k)$, $k = 1, 2, \ldots, K$, gdzie $\tau_k$ jest czasem UTC
odczytanym z komunikatu \texttt{RMC}, a $t_k$ -- wskazaniem zegara
systemowego w chwili odczytu tego komunikatu. Każda para daje jeden pomiar
przesunięcia:
\begin{equation}
  \delta_k = \tau_k - t_k, \qquad k = 1, 2, \ldots, K.
  \label{eq:offset_pomiar}
\end{equation}
Pojedyncze pomiary różnią się nieco od siebie, dlatego jako wynik przyjmuje
się ich średnią arytmetyczną, a rozrzut pomiarów opisuje się odchyleniem
standardowym:
\begin{equation}
  \hat{\delta} = \frac{1}{K} \sum_{k=1}^{K} \delta_k,
  \qquad
  \sigma_{\delta} = \sqrt{\frac{1}{K-1} \sum_{k=1}^{K}
                    \bigl(\delta_k - \hat{\delta}\bigr)^2},
  \label{eq:offset_srednia}
\end{equation}
\begin{objasnienia}
  \item[$\hat{\delta}$] -- wyznaczone przesunięcie zegara w sekundach,
  \item[$\sigma_{\delta}$] -- odchylenie standardowe pomiarów przesunięcia w sekundach,
  \item[$K$] -- liczba pomiarów (przy jednym pomiarze na sekundę sesja trwająca \SI{10}{min} daje ok.~600 pomiarów).
\end{objasnienia}
Poprawiony czas UTC dowolnej zapisanej wiadomości o znaczniku $t$ wynosi
\begin{equation}
  \hat{\tau} = t + \hat{\delta}.
  \label{eq:korekta_czasu}
\end{equation}

Pomiar \eqref{eq:offset_pomiar} nie jest idealny, ponieważ znacznik $t_k$
jest nadawany dopiero po przesłaniu komunikatu łączem szeregowym i odczytaniu
go przez program. Przy prędkości \num{9600}~bit/s przesłanie jednego znaku
(10~bitów) trwa ok.\ \SI{1.04}{ms}, a komunikat NMEA ma do 82 znaków, więc
sama transmisja zajmuje do ok.\ \SI{85}{ms}. Program odczytuje port co
\SI{100}{ms}, co dodaje średnio kolejne \SI{50}{ms}. Wyznaczone przesunięcie
$\hat{\delta}$ jest więc obarczone błędem rzędu \SI{0.1}{s}, co -- jak
pokazano w podrozdziale~\ref{sec:model_przyklad} -- jest wartością
akceptowalną.

\subsection{Model odległości}
\label{sec:model_odleglosc}

\paragraph{Dane wejściowe.}
Z komunikatu SBS typu \texttt{MSG,3} odczytuje się szerokość $\varphi_i$
i długość $\lambda_i$ geograficzną samolotu (w stopniach) oraz jego wysokość
w stopach. Z tematu \texttt{/gps/fix} odczytuje się szerokość
$\varphi_{\mathrm{s}}$, długość $\lambda_{\mathrm{s}}$ i wysokość
$h_{\mathrm{s}}$ stacji (w metrach). Indeks $i$ oznacza samolot, a indeks
$\mathrm{s}$ -- stację.

\paragraph{Przeliczenie jednostek.}
W lotnictwie wysokość podaje się w stopach, a prędkość w węzłach. Na
jednostki układu SI przelicza się je wzorami
\begin{equation}
  h_i = \num{0.3048}\, h_{\mathrm{ft}}, \qquad
  v = \num{0.514444}\, v_{\mathrm{kt}},
  \label{eq:jednostki}
\end{equation}
\begin{objasnienia}
  \item[$h_{\mathrm{ft}}$] -- wysokość samolotu w stopach,
  \item[$h_i$] -- wysokość samolotu w metrach,
  \item[$v_{\mathrm{kt}}$] -- prędkość w węzłach,
  \item[$v$] -- prędkość w metrach na sekundę.
\end{objasnienia}

\paragraph{Odległość po powierzchni Ziemi.}
Ziemię przybliżono kulą o promieniu $R_{\mathrm{Z}} = \SI{6371}{km}$.
Odległość między stacją a punktem na powierzchni Ziemi leżącym dokładnie pod
samolotem oblicza się wzorem haversine:
\begin{equation}
  s = 2 R_{\mathrm{Z}} \arcsin \sqrt{
        \sin^2\frac{\Delta\varphi}{2}
        + \cos\varphi_{\mathrm{s}} \cos\varphi_i
          \sin^2\frac{\Delta\lambda}{2}},
  \label{eq:haversine}
\end{equation}
\begin{objasnienia}
  \item[$s$] -- odległość mierzona po powierzchni Ziemi w metrach,
  \item[$\Delta\varphi$] -- różnica szerokości geograficznych, $\Delta\varphi = \varphi_i - \varphi_{\mathrm{s}}$, w radianach,
  \item[$\Delta\lambda$] -- różnica długości geograficznych, $\Delta\lambda = \lambda_i - \lambda_{\mathrm{s}}$, w radianach,
  \item[$R_{\mathrm{Z}}$] -- średni promień Ziemi w metrach.
\end{objasnienia}
Kąty podawane w stopniach należy przed użyciem we wzorze
\eqref{eq:haversine} zamienić na radiany, mnożąc je przez $\pi/180$.

\paragraph{Odległość skośna.}
Samolot znajduje się na wysokości $h_i$, a stacja na wysokości
$h_{\mathrm{s}}$. Odległość w linii prostej od stacji do samolotu
(odległość skośną) przybliża się, traktując odległość $s$ i różnicę wysokości
jako przyprostokątne trójkąta prostokątnego:
\begin{equation}
  \rho = \sqrt{s^2 + \Delta h^2}, \qquad \Delta h = h_i - h_{\mathrm{s}},
  \label{eq:odleglosc_skosna}
\end{equation}
gdzie $\rho$ jest odległością skośną, a $\Delta h$ -- różnicą wysokości
samolotu i stacji, obie w metrach. Wzory \eqref{eq:jednostki}%
--\eqref{eq:odleglosc_skosna} są zaimplementowane w węźle wyświetlacza
stacji, który na ich podstawie pokazuje odległość do każdego samolotu.

\subsection{Uproszczenia modelu}
\label{sec:model_uproszczenia}

W modelu przyjęto następujące uproszczenia:
\begin{itemize}[nosep]
  \item Ziemia jest kulą, a nie elipsoidą -- powoduje to błąd odległości
        rzędu kilku dziesiątych procenta,
  \item we wzorze \eqref{eq:odleglosc_skosna} pominięto krzywiznę Ziemi;
        błąd rośnie wraz z odległością samolotu od stacji,
  \item wysokość barometryczna z ADS-B jest traktowana jak wysokość nad
        poziomem morza; przy ciśnieniu innym niż standardowe
        (\SI{1013.25}{hPa}) różnica może sięgać kilkudziesięciu metrów,
  \item przesunięcie zegara $\delta$ jest stałe w trakcie sesji (pominięto
        powolny dryf zegara),
  \item stacja jest nieruchoma w trakcie sesji,
  \item pominięto opóźnienie między odebraniem komunikatu ADS-B a nadaniem mu
        znacznika czasu (pojedyncze milisekundy).
\end{itemize}
Uproszczenia te są wystarczające do podglądu na żywo i do analizy zasięgu
stacji. Przy dalszym rozwoju systemu model odległości można zastąpić
dokładniejszym przeliczeniem współrzędnych do lokalnego układu kartezjańskiego
związanego ze stacją, które uwzględnia kształt Ziemi i pozwala dodatkowo
wyznaczyć kierunek obserwacji samolotu.

\subsection{Przykład obliczeniowy}
\label{sec:model_przyklad}

Przyjęto stację w punkcie $\varphi_{\mathrm{s}} = \SI{50.0}{\degree}$,
$\lambda_{\mathrm{s}} = \SI{20.0}{\degree}$, $h_{\mathrm{s}} = \SI{250}{m}$
oraz samolot w punkcie $\varphi_i = \SI{50.3169}{\degree}$,
$\lambda_i = \SI{20.49795}{\degree}$, lecący na wysokości
$h_{\mathrm{ft}} = \num{32800}$~stóp.

Z wzoru \eqref{eq:jednostki} wysokość samolotu wynosi
\begin{equation}
  h_i = \num{0.3048} \cdot \num{32800} \approx \SI{9997.4}{m},
  \qquad
  \Delta h = \num{9997.4} - 250 = \SI{9747.4}{m}.
\end{equation}
Różnice współrzędnych wynoszą $\Delta\varphi = \SI{0.3169}{\degree}$
i $\Delta\lambda = \SI{0.49795}{\degree}$. Wyrażenie pod pierwiastkiem we
wzorze \eqref{eq:haversine} ma wartość ok.\ \num{1.540e-5}, skąd
\begin{equation}
  s \approx \SI{50.00}{km}, \qquad
  \rho = \sqrt{\num{50000}^2 + \num{9747.4}^2} \approx \SI{50.94}{km}.
\end{equation}

Wpływ błędu czasu pokazuje następujące oszacowanie: samolot lecący
z prędkością $v = \SI{230}{m/s}$ w czasie $\Delta\tau$ pokonuje drogę
\begin{equation}
  \Delta x = v \, \Delta\tau .
  \label{eq:blad_czasu}
\end{equation}
Jeżeli zegar stacji myli się o \SI{10}{s} i nie zostanie poprawiony wzorem
\eqref{eq:korekta_czasu}, pozycja przypisana danej chwili będzie błędna
o ok.\ \SI{2.3}{km}. Po korekcji błąd czasu spada do ok.\ \SI{0.1}{s},
a błąd pozycji -- do ok.\ \SI{23}{m}. Uzasadnia to rejestrowanie czasu GPS
razem z danymi ADS-B.

\subsection{Weryfikacja modelu}
\label{sec:model_weryfikacja}

Poprawność modelu odległości można sprawdzić na prostych przypadkach:
dla samolotu położonego dokładnie nad stacją powinno zachodzić $s = 0$
i $\rho = \Delta h$, a dla punktów odległych o \SI{1}{\degree} szerokości
geograficznej $s \approx \SI{111.2}{km}$. Poprawność modelu czasu można
sprawdzić, porównując $\hat{\delta}$ wyznaczone w sesji z dostępem do
internetu (gdy zegar jest synchronizowany przez NTP i $\hat{\delta}$ powinno
być bliskie zera) z sesją bez dostępu do sieci. Szczegółową ocenę na
podstawie zarejestrowanych danych przedstawiono w~rozdziale~\ref{chap:badania}.